\documentclass[a4paper,12pt]{article}
\usepackage[utf8]{inputenc}
\usepackage{amsmath,amssymb,amsfonts,amsthm}
\usepackage{mathrsfs}
\usepackage{enumitem}
\usepackage{multicol}
\usepackage{titlesec}
\usepackage{hyperref}
\usepackage[margin=2.5cm]{geometry}
\setlength{\parindent}{0pt}

\newtheorem{oplossing}{Oplossing}

\begin{document}


\begin{center}
    \Large\textbf{EXAMEN LINEAIRE ALGEBRA} \\[0.5cm]
    \large\textsc{MAANDAG 5 JANUARI 2026} \\[0.5cm]
    \fbox{\textsc{EERSTE BACHELOR FYSICA \& STERRENKUNDE}}
\end{center}


\textbf{Opgave 1.} (6 punten) In deze opgave polsen we naar het inzicht in de leerstof door middel van enkele specifieke vragen.
\begin{enumerate}
    \item[(i)] Zij $V$ een vectorruimte en zij $S \subseteq V$ een lineair onafhankelijke deelverzameling.
    Leg uit waarom de dimensie van $\text{span}(S)$ precies gelijk is aan $|S|$.
    \begin{oplossing}
        Doordat \(S\) lineair onafhankelijk is en doordat \(\operatorname{span}(S)\) per definitie voortbrengend is voor \(S\),
        volgt dat \(S\) een basis is voor \(\operatorname{span}(S)\). Nu is de dimensie van de \(\operatorname{span}(S) = \dim (S) = |S|\).\qed
    \end{oplossing}
    \item[(ii)] In het bewijs van Stelling 3.3.1(i) staat ``$\dots \sum \lambda_{b}f(b)=0$. Dit impliceert dat $\lambda_{b}=0$ voor alle $b \in \mathcal{B}$.
    '' Verklaar dit. Leg hierbij ook uit waarom de assumptie dat de $f(b)$'s twee aan twee verschillend zijn relevant is.
    \begin{oplossing}
       Doordat \(0= \sum_{b \in \mathcal{B}} 0 f(b)\), geldt is wegens de lineair onafhankelijkheid van de \(f(b)\)'s, dat voor ieder andere combinatie \( \sum_{b \in \mathcal{B}}  \lambda_b f(b) = 0\), \(\lambda_b = 0\).
       
       Stel dat de $f(b)$'s niet twee aan twee  verschillend zijn, en \(f(b_1) = f(b_2)\), dan geldt \(\lambda_{b_1}b_1 + \lambda_{b_2} b_2 = 0\) voor \(\lambda_{b_1} = - \lambda_{b_2}\) waarbij ze niet persee \(0\) hoeeven te zijn!
       Hieruit volgt dat het twee aan twee verschillend zijn een noodzakelijke voorwaarde is.
       \qed
    \end{oplossing}
    \item[(iii)] Beschouw twee basissen $\mathcal{B}$ en $\mathcal{B}'$ van een eindig-dimensionale vectorruimte $V$, en zij $Q$ de transitiematrix van $\mathcal{B}$ naar $\mathcal{B}'$ (zie Definitie 4.2.1). Verklaar waarom $Q = A_{\text{id},\mathcal{B}',\mathcal{B}}$ (zie Definitie 4.1.2). Hoe volgt hieruit onmiddellijk dat $Q$ inverteerbaar is?
    \begin{oplossing}
        Neem lemma 4.1.4 (ii) met basis \(\mathcal{B}' = \{b'_1,\dots,b'_n\}\) dan is:
        \[
        \begin{pmatrix}
            \mu_1 \\ \vdots \\ \mu_m
        \end{pmatrix}
        =
        A_{f,\mathcal{B}',\mathcal{C}}
        \begin{pmatrix}
            \lambda'_1 \\ \vdots \\ \lambda'_n
        \end{pmatrix}.
        \]

        Neem nu basis \(\mathcal{B} = \{b_1,\dots, b_n\} \) in plaats van \(\mathcal{C}\), dan verkrijgen we:

        \[
        \begin{pmatrix}
            \lambda_1 \\ \vdots \\ \lambda_n
        \end{pmatrix}
        =
        A_{f,\mathcal{B}',\mathcal{B}}
        \begin{pmatrix}
            \lambda'_1 \\ \vdots \\ \lambda'_n
        \end{pmatrix}.
        \]
        Merk op dat deze basissen dezelfde dimensies hebben en wegens gevolg 3.3.5 isomorf zijn.

        Passen we definitie 4.1.2 hierop toe, dan is \(f(b'_j) = \sum_{i=1}^{n} a_{ij}b_i\).
        \(f(b') = b', \,\forall b' \in \mathcal{B}' \iff f = \operatorname{id} \), waarbij we bij definitie 4.2.1 de transitiematrix toekomen.
        De laatste \(\iff\) relatie geldt uit lemma 3.2.6 (i), \(f\) is volledig bepaald door de beelden van de basiselementen.
        
        Inverteerbaarheid volgt door hetzelfde te doen maar met \(\mathcal{B} \to \mathcal{B}'\) en \(\mathcal{B}' \to \mathcal{B}\).
        \qed
    \end{oplossing} 
    \item[(iv)] Het bewijs van Gevolg 5.3.3 begint met ``Aangezien $\text{rk}(A)=n$ moet ook $\text{rk}(A|w)=n$''. Verklaar dit.
    \begin{oplossing}
        \(A|w\) is een matrix met \(n\) rijen en \(n+1\) kolommen. Uit stelling 5.1.2 kan \(\dim\operatorname{span}(A_1, \dots ,A_n,w ) \leq n\).
        We weten uit lemma 5.1.6 (i) dat \(\operatorname{rk}(A) = \dim(\operatorname{im}) = \dim\operatorname{span}(A_1, \dots ,A_n ) = n \).
        Hieruit volgt dat \(\dim \operatorname{span}(A_1, \dots ,A_n ) \le \dim\operatorname{span}(A_1, \dots ,A_n,w ) \le \dim(K^n) \).
        Dus \(n \le \dim\operatorname{span}(A_1, \dots ,A_n,w ) \le n \implies \dim \operatorname{span}(A_1, \dots ,A_n,w ) = n \implies \operatorname{rk}(A|w) = n\).
        \qed
    \end{oplossing}
    \item[(v)] In het bewijs van Lemma 6.1.6 staat ``Vermits $x$ enkel voorkomt in de elementen op de diagonaal van $B$, is er maar één term van graad $\geq n-1$''. Verklaar dit.
    \begin{oplossing}
        Uit 
        \[\det (Ix-A) = \sum_{\sigma \in S_n} \operatorname{sgn} (\sigma)((x\delta_{\sigma (1),1}-a)_{\sigma (1),1 }\dots , (x\delta_{\sigma (n),n}-a)_{\sigma (n),n})\]
        zien we dat als we permutaties uitvoeren minstens 2 elementen van plaats wisselen.
        Doordat alle \(x\) zich op de diagonaal bevinden is door permutaties uit te voeren dat niet de identiteit is,
        dat de hooste graad \(n-2\) is \(implies\) de \(n-1\) graad wordt bekomen uit  \(\prod_i (x-a_{ii})\).
        \qed
    \end{oplossing}
    \item[(vi)] Het bewijs van Stelling 6.3.4(iii) eindigt met ``en duidelijkerwijze is nu $V=V_{1}\oplus V_{2}\oplus \dots \oplus V_{l}$''. Verklaar dit.
    \begin{oplossing}
        \begin{align*}
            V_1 &= \langle v_1 , \dots, v_{n_1}\rangle \\
            V_2 &= \langle v_{n_1 + 1} , \dots, v_{n_1 + n_2}\rangle \\
            &\vdots \\
            V_i &= \langle v_{n_1 + \dots + n_{i-1} +1} , \dots, v_{n_1 + \dots + n_i}\rangle \\
            V_{i+1} &= \langle v_{n_1 + \dots + n_{i} +1} , \dots, v_{n_1 + \dots + n_{i+1}}\rangle \\
            &\vdots \\
            V_{l} &= \langle v_{n_1 + \dots + n_{l-1} + 1} , \dots, v_{n_1 + \dots + n_{l}}\rangle \\
        \end{align*}
        met \(v_{n_1 + \dots + n_{l}} = v_n\).

        Wegens dat \(\mathcal{B} = \{v_1, \dots , v_n\}\) een basis van \(V\) voorstelt, is \(V_1 + \dots + V_l = V\) en zijn de eigenvectoren lineair onafhankelijk.
        
        Stellen we \(\mathcal{B}_i\) een basis voor \(V_i\), dan is \(\mathcal{B} = \mathcal{B}_1 \cup \dots \cup \mathcal{B}_l\),
        en uit het lineair onafhankelijk zijn \(\mathcal{B}_i \cap \mathcal{B}_j = \emptyset\) voor \((i\neq j )\le l\).
        Uit stelling 2.5.5 volgt dat \(V = V_{1}\oplus V_{2}\oplus \dots \oplus V_{l}\).
        
        \qed
    \end{oplossing}
\end{enumerate}

\vspace{0.5cm}

\textbf{Opgave 2.} (6 punten) Zijn volgende uitspraken waar of vals? Verklaar. (Een antwoord zonder verklaring levert geen punten op. Anderzijds kan een fout antwoord met een zinvolle deelredenering wél punten opleveren.) \\
\textit{Opmerking: een uitspraak is pas waar als ze altijd waar is. Omgekeerd geldt dus dat een uitspraak vals is zodra er een voorbeeld bestaat waarvoor ze vals is; we spreken dan van een tegenvoorbeeld.}
\begin{enumerate}
    \item[(i)] Als $A$ en $B$ vierkante matrices zijn, dan is de rang van $AB$ gelijk aan de rang van $BA$.
    \begin{oplossing}
        \textbf{Fout!}
        Stel

        \(A =
        \begin{pmatrix}
            1 & 1 \\
            2 & 2
        \end{pmatrix}
        \)
        en 
        \(B =
        \begin{pmatrix}
            1 & 0 \\
            -1 & 0
        \end{pmatrix}
        \).

        Dan is 
        \(AB =
        \begin{pmatrix}
            0 & 0 \\
            0 & 0
        \end{pmatrix}
        \)
        en
        \(BA =
        \begin{pmatrix}
            1 & 1 \\
            -1 & -1
        \end{pmatrix}
        \).

        Waarbij \(\operatorname{rk}(AB) = 0\) en \(\operatorname{rk}(BA) = 1\).\qed
    \end{oplossing}
    \item[(ii)] Zij $V$ en $W$ twee eindig-dimensionale vectorruimten. Als $f:V \to W$ en $g:W \to V$ twee lineaire afbeeldingen zijn zodanig dat $g \circ f = \text{id}_{V}$, dan is $\dim V = \dim W$.
    \begin{oplossing}
        \textbf{Fout!}
        Stel \(f: \begin{pmatrix}
            a \\
            b
        \end{pmatrix}
        \mapsto
        \begin{pmatrix}
            a \\ b \\ a+b
        \end{pmatrix}
        \)
        en
        \(g: \begin{pmatrix}
            a \\ b \\ c 
        \end{pmatrix}
        \mapsto 
        \begin{pmatrix}
            a \\ b
        \end{pmatrix}\).

        Dan geldt:
        \(g\left( f\left(\begin{pmatrix}
            a \\ b
        \end{pmatrix}\right) \right)
        =
        g\left( \begin{pmatrix}
            a \\ b \\ a+b
        \end{pmatrix}\right) 
        =
        \begin{pmatrix}
            a \\ b
        \end{pmatrix}
        = \operatorname{id}_V.
        \)

        Maar \(\dim V \neq \dim W\). \qed
    \end{oplossing}
    \item[(iii)] Zij $f:\mathbb{R}^{3} \to \mathbb{R}^{3}$ een lineaire afbeelding met $\det(f)=0$. Als er twee lineair onafhankelijke vectoren $v,w \in \mathbb{R}^{3}$ bestaan zodat $f(v)=w$ en $f(w)=v$, dan is $f$ diagonaliseerbaar.
    \begin{oplossing}
        \textbf{Juist!}

        \(\det f = 0 \implies \lambda_1 = 0\).

        \(f(w+v) = w+v \implies \lambda_2 = 1\).

        \(f(w-v) = -(w-v) \implies \lambda_3 = -1\).

        Drie verschillende eigenwaarden voor \(\mathbb{R}^3 \implies f\) diagonaliseerbaar. \qed 
    \end{oplossing}
    \item[(iv)] Beschouw een eindig-dimensionale reële inproductruimte $V$ met deelruimten $U$ en $W$. Als $V=U \oplus W$, dan is ook $V=U^{\perp} \oplus W^{\perp}$.
    \begin{oplossing}
        \textbf{Waar!}
        Via dimensies:

        \(\dim W^\perp = \dim V - \dim W\)

        \(\dim U^\perp = \dim V - \dim U\)

        \(\dim V = \dim U + \dim W \)
        
        \(\implies \dim V = \dim V - \dim U^\perp + \dim V - \dim W^\perp \implies \dim V = \dim U^\perp + \dim W^\perp\).

        Nu tonen we aan dat \(U^\perp \cap W^\perp = \{0\}\).

        Stel \(t \in U^\perp \cap W^\perp \), dan \(\langle t,w \rangle = \langle t,u \rangle = 0 ,\, \forall w \in W, \, u \in U \implies \langle t,v \rangle = 0 , \, \forall v \in V .\)
    
        Dit laatste kan enkel als \(v=0\). \qed
    \end{oplossing}
\
\vspace{0.5cm}
\textbf{Begin een nieuw geruit blad.}
\vspace{0.5cm}

\textbf{Opgave 3.} (4 punten) Zij $V$ en $W$ vectorruimten over $\mathbb{R}$, met respectieve basissen $\mathcal{B}=\{b_{1},b_{2},b_{3}\}$ en $\mathcal{C}=\{c_{1},c_{2}\}$. Zij $f:V \to W$ een lineaire afbeelding met matrixvoorstelling ten opzichte van $\mathcal{B}$ en $\mathcal{C}$:
\[
A_{f} = \begin{pmatrix}
1 & 0 & 1 \\
0 & 2 & -1
\end{pmatrix}
\]
\begin{enumerate}
    \item[(i)] Bepaal de kern en het beeld van $f$ (uitgedrukt ten opzichte van $\mathcal{B}$ en $\mathcal{C}$).
    \item[(ii)] We definiëren
    \[
    v_{1} := b_{1} - 2b_{2}, \quad v_{2} := b_{1} + b_{2} + b_{3}, \quad v_{3} := b_{3} - b_{1}
    \]
    en
    \[
    w_{1} := \frac{1}{2}(c_{1} + c_{2}), \quad w_{2} := \frac{1}{2}(c_{2} - c_{1}).
    \]
    Toon aan dat $\mathcal{D}:=\{v_{1},v_{2},v_{3}\}$ en $\mathcal{E}:=\{w_{1},w_{2}\}$ basissen zijn voor respectievelijk $V$ en $W$.
    \item[(iii)] Bepaal de matrixvoorstelling $B_{f}$ van $f$ ten opzichte van $\mathcal{D}$ en $\mathcal{E}$.
\end{enumerate}

\vspace{0.5cm}

\textbf{Opgave 4.} (4 punten) Beschouw volgende matrix over $\mathbb{C}$:
\[
A = \begin{pmatrix}
2i+3 & 0 & -i-1 \\
0 & 1 & 0 \\
6i+6 & 0 & -3i-2
\end{pmatrix}
\]
\begin{enumerate}
    \item[(i)] Toon aan dat $\chi_{A}(x) = (x-1)^{2}(x+i)$.
    \item[(ii)] Toon aan dat $A$ diagonaliseerbaar is, bepaal de eigenruimten van $A$ en bepaal een diagonaalmatrix $D$ en een transitiematrix $P$ zodanig dat $A = PDP^{-1}$.
    \item[(iii)] Leid hieruit af dat $A$ inverteerbaar is en bepaal $A^{-1}$ zonder gebruik te maken van rijreductie of de adjunct van $A$.
    \item[(iv)] Bepaal alle $n \in \mathbb{Z}$ zodat $A^{n} = I_{3}$.
\end{enumerate}



\begin{center}
    \Large\textbf{EXAMEN LINEAIRE ALGEBRA} \\[0.5cm]
    \large\textsc{Maandag 6 januari 2025} \\[0.5cm]
    \fbox{\textsc{Eerste bachelor Fysica}}
\end{center}

\vspace{1cm}

\begin{center}
    \textsc{Theorie}
\end{center}

\textbf{Opgave 1.} (4 punten) In deze opgave wordt gevraagd om een aantal kleine uitbreidingen op de cursusnota's te bewijzen.
\begin{itemize}
    \item[(i)] Toon de volgende bewering aan: Stel $\{b_1, \dots, b_n\}$ is een basis van een vectorruimte $V$ en $\{c_1, \dots, c_n\}$ een verzameling van vectoren in een vectorruimte $W$. Dan bestaat er een unieke lineaire afbeelding $f : V \to W$ zodat $f(b_i) = c_i$ voor alle $1 \leq i \leq n$.
    \begin{oplossing}
    Lemma 3.2.6(i): $f$ is volledig bepaald door de beelden v.d. basis $f(b) \quad \forall b \in B$.

    

    $v \in V: v = \sum_i^n \lambda_i b_i \implies f(v) = \sum_i^n \lambda_i f(b_i) \quad (f \text{ is lineair})$


    doordat $v$ willekeurig is, is $f$ volledig bepaald.
    Uniciteit volgt uit dat alle $v \in V$ op een unieke manier als lineaire combinatie uit te schrijven zijn. 
    \end{oplossing} \qed

    \item[(ii)] Stel $\mathcal{B}$ is een basis van een vectorruimte $V$ en $f : V \to V$ een endomorfisme zodat de geassocieerde matrix $A_{f,\mathcal{B}}$ gelijk is aan
    \[
    \begin{pmatrix}
    \lambda & 1 & 0 & 0 \\
    0 & \lambda & 1 & 0 \\
    0 & 0 & \lambda & 1 \\
    0 & 0 & 0 & \lambda
    \end{pmatrix}.
    \]
    Preciseer en verklaar voor welke basis $\mathcal{C}$ van $V$ de geassocieerde matrix $A_{f,\mathcal{C}}$ gegeven wordt door
    \[
    \begin{pmatrix}
    \lambda & 0 & 0 & 0 \\
    1 & \lambda & 0 & 0 \\
    0 & 1 & \lambda & 0 \\
    0 & 0 & 1 & \lambda
    \end{pmatrix}.
    \]
    Stel eveneens de overgangsmatrix van $\mathcal{B}$ naar $\mathcal{C}$ op.
    
    \item[(iii)] Zij $A \in M_n(K)$ en $T_A$ de verzameling van matrices $B$ in $M_n(K)$ zodat $ABA^{-1}$ een diagonaalmatrix is. Toon aan dat $T_A$ een $n$-dimensionale deelruimte is van $M_n(K)$.
    \begin{oplossing}
        Hier zit ik nog vast!
    \end{oplossing}
    \begin{oplossing}
        \(T_A\) is een deelruimte als: 
        \[
        \lambda B_1 + \mu B_2 \in T_A \qquad \forall \mu, \lambda \in K \text{ en } B_1 , B_2 \in T_A.
        \]
        We heten de diagonaalmatrix van een matrix \(B_i \in T_A: \operatorname{diag}(B_i)\). 
        De notatie is vrij vanzichzelf uitsprekend.

        Dan is
        \begin{align*}
            A \left(\lambda B_1 + \mu B_2\right) A^{-1} 
            &= A \lambda B_1 A^{-1} + A \mu B_2 A^{-1} \\
            &=  \lambda A B_1 A^{-1} + \mu A  B_2 A^{-1} \\
            &=   \lambda \operatorname{diag}(B_1) + \mu \operatorname{diag}(B_2) \\
            &=  \operatorname{diag}(\lambda B_1) + \mu \operatorname{diag}( \mu B_2) \\
            &= \operatorname{diag}(\lambda B_1 + \mu B_2) \in T_A.
        \end{align*}\qed
    \end{oplossing}
    \item[(iv)] Zij $V$ een $n$-dimensionale vectorruimte. Toon aan dat de doorsnede van $k < n$ deelruimten van dimensie $n-1$ minstens dimensie $n-k$ heeft.
    
    \begin{oplossing}
        Idee: bewijzen via inductie \dots

        Noteer de deelruimten \(W_i\) met \(\dim (W_i) = n-1\).

        We beginnen bij \(k = 2\). Neem \(W_1 \le V \), \(W_2 \le V\) met \(\dim(W_1) = \dim(W_2) = n-1\).
        Uit de dimensiestelling voor deelruimten (2.5.6) weten we dat:
        \[
        \dim(W_1 +W_2) = \dim(W_1) + \dim (W_2) - \dim(W_1 \cap W_2).
        \]
        Opmerking: \(\dim(W_1 + W_2) \leq \dim(V) = n \) wegens \(W_i\) deelruimten zijn van \(V\).
        Hieruit volgt
        \begin{align*}
            \dim(W_1 \cap W_2)
            &= \dim(W_1) + \dim (W_2) - \dim(W_1 +W_2) && [- \dim(W_1 + W_2)\geq - \dim(V)] \\
            &\geq \dim(W_1) + \dim (W_2) - \dim(V) \\
            &\geq 2(n-1) -n = n-2.
        \end{align*}

        Neem nu \(W'_s\), de doorsnede van \(W_1 \cap W_2 \cap \dots \cap W_s \) met \(s < k < n\) deelruimten van dimensie \(n-1\), en stel \(\dim(W'_s) \geq n-s\).
        Dan tonen we aan dat \(\dim(W'_s \cap W_{s+1}) \geq n-(s+1)\):
        \begin{align*}
            \dim(W'_s \cap W_{s+1}) 
            &= \dim(W'_s) + \dim (W_{s+1}) - \dim(W'_s + W_{s+1})\\
            &\geq \dim(W'_s) + \dim (W_{s+1}) - \dim(V)\\
            &\geq (n-s) + (n-1) - n = n- (s+1).
        \end{align*}
    \end{oplossing}
    Bijgevolg is \(\dim\left(W_1 \cap W_2 \cap \dots \cap W_s \cap W_{s+1} \cap \dots \cap W_k \right) \leq n-k \), \((k < n)\).
    \qed
\end{itemize}

\vspace{0.5cm}

\textbf{Opgave 2.} (2 punten) In deze vraag tonen we aan dat voor diagonaliseerbare matrices $A, B \in M_n(\mathbb{C})$ de lineaire afbeeldingen $L_A$ en $L_B$ tegelijk diagonaliseerbaar zijn t.o.v.\ eenzelfde basis als en slechts als $AB = BA$.
\begin{itemize}
    \item[(i)] Zij $\lambda$ een eigenwaarde van $A$. Toon aan dat als $AB = BA$, dan is $L_B(\text{Eigr}_\lambda(A)) \subseteq \text{Eigr}_\lambda(A)$.
    \begin{oplossing}
        \[A v = \lambda v \implies BAv = B\lambda v = \lambda (Bv). \]
        \[\implies A(Bv) = \lambda(Bv).\]
        \(Bv\) is dus ook een eigenvector van \(A\) met dezelfde eigenwaarde \(\lambda\).
        \begin{align*}
            L_B (\operatorname{Eigr}_\lambda (A)) 
            &= \{ v\in K^n | B(\lambda I_n - A)v = 0)\} \\
            &= \{ v\in K^n | (B\lambda I_n - BA)v = 0)\} \\
            &=\{ v\in K^n | (\lambda I_n B - A B)v = 0)\} \\
            &= \{ (Bv)\in K^n | (\lambda I_n - A)(Bv) = 0)\}\\
            &\subseteq \operatorname{Eigr}_\lambda (A).
        \end{align*}\qed
    \end{oplossing}
    \item[(ii)] Toon aan dat als $AB = BA$, dan bestaat er een basis van $\mathbb{C}^n$ zodat de geassocieerde matrices van de lineaire afbeeldingen $L_A$ en $L_B$ diagonaalmatrices zijn. Gebruik hiervoor het vorige puntje en het volgende lemma (dat je mag gebruiken zonder bewijs) ``Zij $W \leq \mathbb{C}^n$ een deelruimte en $f$ een endomorfisme van $V$ zodat $f|_W$ een endomorfisme van $W$ is. Dan is $f|_W$ diagonaliseerbaar, als $f$ diagonaliseerbaar is.''
    \begin{oplossing}
        \textbf{Nog te doen!}
    \end{oplossing}
    \item[(iii)] Toon aan dat als $L_A$ en $L_B$ diagonaliseerbaar zijn t.o.v.\ dezelfde basis voor $\mathbb{C}^n$, dan is $AB = BA$.
    \begin{oplossing}
        Stel \(P^{-1} A P = \operatorname{diag}(\lambda_1,\dots,\lambda_n)\) en \(P^{-1} B P = \operatorname{diag}(\lambda'_1,\dots,\lambda'_n)\) 
        Dan:
        \begin{align*}
            (P^{-1} A P)(P^{-1} B P) 
            &= P^{-1} AB P\\
            &=  \operatorname{diag}(\lambda_1,\dots,\lambda_n)  \operatorname{diag}(\lambda'_1,\dots,\lambda'_n)\\
            &= \operatorname{diag}(\lambda_1\lambda'_1,\dots,\lambda_n\lambda'_n),\\
            \\
            (P^{-1} B P)(P^{-1} A P) 
            &= P^{-1} BA P\\
            &=  \operatorname{diag}(\lambda'_1,\dots,\lambda'_n)  \operatorname{diag}(\lambda_1,\dots,\lambda_n)\\
            &= \operatorname{diag}(\lambda_1\lambda'_1,\dots,\lambda_n\lambda'_n). \\
            \\
            \implies AB=BA.
        \end{align*}
        Waarbij we gebruik hebben gemaakt van eigenschap (K8), de commutativiteit van de scalaire vermenigvuldiging. \qed
    \end{oplossing}
\end{itemize}

\vspace{0.5cm}

\textbf{Opgave 3.} (4 punten) Zijn volgende uitspraken juist of fout? Indien juist, geef een argument; indien fout, geef een tegenvoorbeeld. (Enkel ``juist'' of ``fout'' antwoorden zonder verklaring levert je geen punten op.)
\begin{itemize}
    \item[(i)] Zij $W_1, W_2, W_3$ deelruimten van een eindigdimensionale vectorruimte $V$ met $W_1 \oplus W_2 = W_1 \oplus W_3$, dan is $W_2 \cong W_3$.
    \begin{oplossing}
        \textbf{Juist!}
        We gebruiken stelling 2.5.4 (ii). Stel \(W \le V\) en \(W = W_1 \oplus W_2 = W_1 \oplus W_3\) met \(\dim(W) = n\).
        Dan:
        \[
        \dim W = \dim W_1 + \dim W_2 = \dim W_1 + \dim W_3 \implies \dim W_2 = \dim W_3.
        \]
        Uit gevolg 3.3.5 volgt het te bewijzen. \qed
    \end{oplossing}
    \item[(ii)] Zij $A, B \in M_n(K)$ met $\lambda_A$ (resp.\ $\lambda_B$) een eigenwaarde van $A$ (resp.\ $B$). Dan is $\lambda_A \lambda_B$ een eigenwaarde van $AB$.
    \begin{oplossing}
        \textbf{Fout!} Stel \(A en B\) hebben elk \(n\) verschillende eigenwaarden. Dan stelt deze uitspraak dat \(AB\), \(n^2\) eigenwaarden heeft.
        Maar \(AB\) heeft hoogstens \(\operatorname{rk}=n\), waardoor de karakteristieke veelterm van \(AB\) niet meer dan \(n\) eigenwaarden kan hebben.
        Wat strijdig is met de uitspraak. 

        Vb:
        \(A = \begin{pmatrix}
            1 & 0 \\
            0 & 5
        \end{pmatrix}
        \)
        en 
        \(B = \begin{pmatrix}
            2 & 0 \\
            0 & 3
        \end{pmatrix}
        \)
        Overduidelijk is:
        \(\lambda_{A,1} = 1, \lambda_{A,2} = 5 \) en \(\lambda_{B,1} = 2, \lambda_{B,2} = 3\).

        Nu is 
        \(AB = \begin{pmatrix}
            2 & 0 \\
            0 & 15
        \end{pmatrix}
        \) 
        met 
        \(\lambda_{AB,1} = 2, \lambda_{AB,2} = 15\). 

        Het is duidelijk dat \(\lambda_{A,1}\lambda_{B,2} = 5 \notin \lambda_{AB}\).\qed
    \end{oplossing}
    \item[(iii)] Zij $A$ een diagonaliseerbare matrix en $B$ een matrix rij-equivalent met $A$, dan is $B$ ook diagonaliseerbaar.
    \begin{oplossing}
        \textbf{Fout!}
        In de cursus is matrix \(B\) in voorbeeld 6.3.5 (2) niet diagonaliseerbaar, met
        \(B = \begin{pmatrix}
            1 & 0 & 0 \\
            0 & -1 & 0\\
            0 & 1 & -1
        \end{pmatrix}\).

        Maar een rij-operatie \(R_3 \to R_3 + R_2\) bekomt volgende matrix:

        \(A = \begin{pmatrix}
            1 & 0 & 0 \\
            0 & -1 & 0\\
            0 & 0 & -1
        \end{pmatrix}\),
        een diagonaalmatrix.

        Wat een strijdig voorbeeld is voor deze uitspraak. \qed
    \end{oplossing}
    \item[(iv)] Zij $V$ een vectorruimte en $f$ een endomorfisme van $V$ zodat voor elke $v \in V$ geldt dat $f(f(v)) = f(v)$. Dan is $\ker f = 0$.
    \begin{oplossing}
        \textbf{Fout!} Een endomorfisme is een lineaire afbeelding, maar niet algemeen bijectief!

        Vb: \(f: V \to V: \begin{pmatrix}
            x \\
            y
        \end{pmatrix}
        \mapsto
        \begin{pmatrix}
            x \\
            0
        \end{pmatrix}
         \).
         
         Dan is \(f \left(\begin{pmatrix}
            x \\
            y
        \end{pmatrix}\right) 
        = \begin{pmatrix}
            x \\
            0
        \end{pmatrix}
         \)
        en \(f\left(f \left(\begin{pmatrix}
            x \\
            y
        \end{pmatrix}\right)\right)
        = \begin{pmatrix}
            x \\
            0
        \end{pmatrix}
         \)

         en \(\ker f = \{\begin{pmatrix}
            0 \\
            y
        \end{pmatrix} | y \in \mathbb{R}\}\).

        Wat opnieuw een tegenvoorbeeld is. \qed
    \end{oplossing}
\end{itemize}

\newpage

\begin{center}
    \textsc{Oefeningen}
\end{center}

\fbox{Begin een nieuw geruit blad.}

\vspace{0.5cm}

\textbf{Opgave 4.} (6 punten) Beschouw de matrix
\[
A = \begin{pmatrix}
2 & 2 & -1 \\
2 & 5 & -2 \\
-1 & -2 & 2
\end{pmatrix}.
\]
Zij $L_A$ de lineaire afbeelding op $\mathbb{R}^3$ gegeven door linksvermenigvuldiging met $A$. Definieer $\langle ., . \rangle : \mathbb{R}^3 \times \mathbb{R}^3 \to \mathbb{R}$ als volgt:
\[
\langle x, y \rangle = x^t A y.
\]
In deel (i) en (iv) wordt aangetoond dat $\langle ., . \rangle$ een inproduct is.
\begin{itemize}
    \item[(i)] Toon aan dat $\langle ., . \rangle$ symmetrisch is, en lineair in het eerste argument.
    \item[(ii)] Toon aan dat $(1, 1, 3)^t$ een eigenvector is van $L_A$ met eigenwaarde $1$.
    \item[(iii)] Toon aan dat alle eigenwaarden van $L_A$ strikt positief zijn.
    \item[(iv)] Toon met behulp van Gevolg 8.4.4 aan dat $\langle x, y \rangle := x^t A y$ positief-definiet is.
    \item[(v)] Bereken, ten opzichte van het gegeven inproduct, de hoek tussen de vectoren $x = (0, 2, 3)^t$ en $y = (-3, 2, 1)^t$.
\end{itemize}

\vspace{0.5cm}

\textbf{Opgave 5.} (4 punten) Zij $V$ de verzameling van afbeeldingen van $\{1, 2, 3, 4\}$ naar $\mathbb{Q}$. Definieer een optelling en scalaire vermenigvuldiging als volgt. Zij $f, g \in V$ en $\lambda \in \mathbb{Q}$, dan
\begin{align*}
(f + g)(i) &= f(i) + g(i), \\
(\lambda * f)(i) &= \lambda * f(i).
\end{align*}
Met deze bewerkingen vormt $V$ een vectorruimte.
\begin{itemize}
    \item[(i)] Ga eigenschappen (V2), (V3) en (V5) van Definitie 2.1.1 na voor $V$.
    \item[(ii)] Zij $\varphi \in V^\ast$. Noteer
    \[
    W = \{f \in V \mid \varphi(f) = f(2)\}.
    \]
    Is $W$ een deelruimte van $V$? Motiveer je antwoord.
    \item[(iii)] Zij $f, g, h \in V$ gegeven door
    \begin{align*}
    f(1) &= 0, & f(2) &= -1, & f(3) &= 3, & f(4) &= 2, \\
    g(1) &= 1, & g(2) &= -3, & g(3) &= 2, & g(4) &= -2, \\
    h(1) &= 2, & h(2) &= 0, & h(3) &= 0, & h(4) &= 1.
    \end{align*}
    Is $\{f, g, h\}$ een lineair onafhankelijke verzameling van $V$? Motiveer je antwoord.
\end{itemize}



\begin{center}
    \Large\textbf{INHAALEXAMEN LINEAIRE ALGEBRA} \\[0.5cm]
    \large\textsc{Vrijdag 12 januari 2024} \\[0.5cm]
    \fbox{\textsc{eerste bachelor Fysica \& Sterrenkunde}}
\end{center}

\vspace{1cm}

\begin{center}
    \textsc{Theorie}
\end{center}

\textbf{Opgave 1.} (3 punten) In deze opgave wordt gevraagd om een aantal argumenten of overgangen uit de cursusnota's in detail te verklaren of verder op te bouwen.
\begin{itemize}
    \item[(i)] In stelling 6.3.5 gebruikt men dat de unie van basissen van eigenruimten een lineair afhankelijke verzameling is. Leg uit waarom en 6.3.4.(1) is niet afdoende.
    \item[(ii)] Verderbouwend op stelling 5.5.2, toon aan dat de dimensie van de grootste minor van een vierkante matrix $A$ met niet-nul determinant gelijk is aan de rang van $A$.
    \item[(iii)] Onderstel $A, B \in M_n(K)$. Toon aan dat als $C = AB$ en $C = (C_1 \dots C_n)$ de kolommen van $C$, dan is $C_k = b_{1,k}A_1 + \dots + b_{n,k}A_n$, waarbij de $A_i$ de kolommen van $A$ zijn.
    \begin{oplossing}
        \(C = AB = \sum_{j=1}^{n} a_{ij}b_{jk}\).

        \(C_k 
        = \left(\sum_{j=1}^{n} a_{ij}b_{j}\right)_k  
        = \sum_{j=1}^{n} b_{j,k}a_{ij}
        = \sum_{j=1}^{n} b_{j,k}A_j\). \qed
    \end{oplossing}
\end{itemize}

\vspace{0.5cm}

\textbf{Opgave 2.} (4 punten) We noemen matrices $A, B \in M_n(K)$ simultaan, als er een $S \in M_n(K)$ bestaat zodat zowel $S^{-1}AS$ en $S^{-1}BS$ diagonaalmatrices zijn.
\begin{itemize}
    \item[(i)] Geef een voorbeeld van 2 diagonaliseerbare matrices die niet simultaan zijn.
    \begin{oplossing}
        Dit zijn 2 diagonaliseerbare matrices die niet commuteren.
        Bv:
        \(
        A = \begin{pmatrix}
            -1 & 1 \\
            0 & 1
        \end{pmatrix}
        \)
        en 
        \(
        B = \begin{pmatrix}
            1 & 0 \\
            1 & 0
        \end{pmatrix}
        \).

        Hier is 
        \(
        AB = \begin{pmatrix}
            0 & 0 \\
            1 & 0
        \end{pmatrix}
        \)
        en
        \(
        BA = \begin{pmatrix}
            -1 & 1 \\
            -1 & 1
        \end{pmatrix}
        \)

        Dus \(AB \neq BA\).
    \end{oplossing}
    \item[(ii)] Toon aan dat als $A$ en $B$ simultaan zijn, dat dan $AB = BA$.
    \begin{oplossing}
        Stel \(S^{-1} A S = \operatorname{diag}(\lambda_1,\dots,\lambda_n)\) en \(S^{-1} B S = \operatorname{diag}(\lambda'_1,\dots,\lambda'_n)\) 
        Dan:
        \begin{align*}
            (S^{-1} A S)(S^{-1} B S) 
            &= S^{-1} AB S\\
            &=  \operatorname{diag}(\lambda_1,\dots,\lambda_n)  \operatorname{diag}(\lambda'_1,\dots,\lambda'_n)\\
            &= \operatorname{diag}(\lambda_1\lambda'_1,\dots,\lambda_n\lambda'_n),\\
            \\
            (S^{-1} B S)(S^{-1} A S) 
            &= S^{-1} BA S\\
            &=  \operatorname{diag}(\lambda'_1,\dots,\lambda'_n)  \operatorname{diag}(\lambda_1,\dots,\lambda_n)\\
            &= \operatorname{diag}(\lambda_1\lambda'_1,\dots,\lambda_n\lambda'_n). \\
            \\
            \implies AB=BA.
        \end{align*}
        Waarbij we gebruik hebben gemaakt van eigenschap (K8), de commutativiteit van de scalaire vermenigvuldiging. \qed
    \end{oplossing}
    \item[(iii)] Stel dat $AB = BA$ en $A$ $n$ verschillende eigenwaarden heeft. Toon aan dat $A$ en $B$ simultaan zijn. Tip: Bekijk wat er gebeurt met een eigenvector van $A$ als je die links vermenigvuldigt met $B$. Wat betekent dit voor de eigenruimten van $A$.
     \begin{oplossing}
        \[A v = \lambda v \implies BAv = B\lambda v = \lambda (Bv). \]
        \[\implies A(Bv) = \lambda(Bv).\]
        \(Bv\) is dus ook een eigenvector van \(A\) met dezelfde eigenwaarde \(\lambda\).
        \[B v = \mu v \implies ABv = A\mu v = \mu (Av). \]
        \[\implies B(Av) = \mu(Av).\]
        Idem is \(Av\) ook een eigenvector van \(B\) met dezelfde eigenwaarde \(\mu\).
        \begin{align*}
            L_B (\operatorname{Eigr}_\lambda (A)) 
            &= \{ v\in K^n | B(\lambda I_n - A)v = 0)\} \\
            &= \{ v\in K^n | (B\lambda I_n - BA)v = 0)\} \\
            &=\{ v\in K^n | (\lambda I_n B - A B)v = 0)\} \\
            &= \{ (Bv)\in K^n | (\lambda I_n - A)(Bv) = 0)\}\\
            &\subseteq \operatorname{Eigr}_\lambda (A).
        \end{align*}
        Als we dit omgekeerd doen zien we dat ook \(L_A (\operatorname{Eigr}_\mu (B))\subseteq \operatorname{Eigr}_\mu (B)\).

        Doordat \(A\) \(n\) verschillende eigenwaarden heeft, zijn de eigenvectoren en dus ook de eigenruimten lineair onafhankelijk en vormen ze een basis.
        
        Dit betekend dat er een basis bestaat zodat \(A\) en \(B\) diagonaal matrices worden. 
        \qed
    \end{oplossing}
\end{itemize}

\vspace{0.5cm}

\textbf{Opgave 3.} (3 punten) Zijn volgende uitspraken juist of fout? Indien juist, geef een argument; indien fout, geef een tegenvoorbeeld. (Enkel ``juist'' of ``fout'' antwoorden zonder verklaring levert je geen punten op.)
\begin{itemize}
    \item[(i)] Als $K$-vectorruimte zijn $K[X]$ en $\bigoplus_{i\in\mathbb{N}}K$ isomorf.
    \begin{oplossing}
        \textbf{Juist} \(f: \bigoplus_{i\in\mathbb{N}}K \to K[X]:  \begin{pmatrix}
            0\\
            1\\
            \vdots \\
            i \in \mathbb{N} \\
            \vdots \\
        \end{pmatrix} \mapsto 
        \begin{pmatrix}
            x^{0}\\
            x^{1} \\
            \vdots \\
            x^{i} \\
            \vdots \\
        \end{pmatrix}
        \).
    \end{oplossing}
    \item[(ii)] Als $A, B, C$ deelruimten zijn van een vectorruimte $V$ en $A \cap B = A + C$ en $B \cap C = A + B$ dan is $A = B = C$.
    \begin{oplossing}
        \textbf{Juist!}

        \(B \subseteq A + B = B \cap C \implies A \subseteq B \subseteq C\).

        \(A \subseteq A + C = A \cap B \implies C \subseteq A \subseteq B\).

        Uit \(C \subseteq A \text{ en } A \subseteq B \subseteq C \text{ is } C = A \).
        
        Uit \( B \subseteq C \text{ en } C \subseteq A \subseteq B \text{ is } C = B\). 
        
        Hieruit volgt $A = B = C$. \qed
    \end{oplossing}
    \item[(iii)] Als $A, B$ diagonaliseerbare $n \times n$-matrices zijn dan is $A+B$ diagonaliseerbaar.
    \begin{oplossing}
        \textbf{Fout!}
        In de cursus is matrix \(B\) (\textbf{die we nu \(A+B\) benoemen}) in voorbeeld 6.3.5 (2) niet diagonaliseerbaar, met
        \(A  = \begin{pmatrix}
            1 & 0 & 0 \\
            0 & -1 & 0\\
            0 & 1 & -1
        \end{pmatrix}\).

        Maar met bv: 
        \(A  = \begin{pmatrix}
            -5 & 0 & 0 \\
            0 & 1 & 0\\
            0 & 1 & 3
        \end{pmatrix}\)
        en
        \( B = \begin{pmatrix}
            6 & 0 & 0 \\
            0 & -2 & 0\\
            0 & 0 & -4
        \end{pmatrix}\),
        zijn \(A\) en \(B\) duidelijk diagonaliseerbaar.
        
        Maar \(A+B\) is niet diagonaliseerbaar, dit is dus een geldig tegenvoorbeeld. \qed
    \end{oplossing}
\end{itemize}

\newpage

\begin{center}
    \textsc{Oefeningen}
\end{center}

\fbox{Begin een nieuw enkel geruit blad.}

\vspace{0.5cm}

\textbf{Opgave 4.} (5 punten) Beschouw volgend stelsel over $\mathbb{Q}$ met parameter $a \in \mathbb{Q}$.
\[
\begin{cases}
(a + 1)x + y + z = a \\
x + (a + 1)y + z = 1 - 2a \\
x + y + (a + 1)z = a - 1
\end{cases}
\]
\begin{itemize}
    \item[(i)] Bespreek de oplosbaarheid en het aantal vrijheidsgraden (voor de onbekenden $x, y$ en $z$) voor alle waarden van de parameter $a$.
    \item[(ii)] Bepaal voor elk van de oplossingsverzamelingen de affiene deelruimte $D$ die erdoor wordt opgespannen. Geef de affiene deelruimte als $v + \text{span}(b_1, \dots, b_k)$, met $v$ een vector en $\{b_1, \dots, b_k\}$ een basis voor de onderliggende vectorruimte van $D$. Geef aan in welke gevallen de affiene deelruimte ook een vectordeelruimte is.
\end{itemize}

\vspace{0.5cm}

\fbox{Begin een nieuw enkel geruit blad.}

\vspace{0.5cm}

\textbf{Opgave 5.} (5 punten) Zij $L_A$ de lineaire afbeelding van $\mathbb{R}^3 \to \mathbb{R}^3$ bekomen door linkse vermenigvuldiging met de matrix
\[
A = \begin{pmatrix}
2 & 0 & 0 \\
0 & 2 & -1 \\
0 & 4 & -3
\end{pmatrix}.
\]
We beschouwen dus de afbeelding
\[
L_A : \mathbb{R}^3 \to \mathbb{R}^3 :
\begin{pmatrix} x \\ y \\ z \end{pmatrix}
\mapsto
\begin{pmatrix}
2 & 0 & 0 \\
0 & 2 & -1 \\
0 & 4 & -3
\end{pmatrix}
\begin{pmatrix} x \\ y \\ z \end{pmatrix}.
\]
Bepaal een nieuwe basis $\mathcal{B}$ voor $\mathbb{R}^3$ zodat $L_A$ ten opzichte van de nieuwe basis wordt voorgesteld door een diagonaalmatrix.

\end{document}